\section{Conclusion}
Complete windowed analysis can advance over an exact hop as one bounded
sequence covering preparation, transform, magnitude processing, smoothing, and
publication. Whole-pipeline accounting removes the boundary bursts and cadence
error of FFT-only scheduling, while input retention and mailbox ownership
protect the analyzed samples and the published spectra.

A descriptive study separates this contract from implementation cost. Matched
scheduling reduces typical hop peaks, and native hybrids and shorter
completion horizons offer different tradeoffs between cost and age. With four
aligned Rack modules on one engine thread, the production pipeline
substantially reduces block concentration relative to native batch execution.
Synchronization, staggering, and release lateness nevertheless reshape
host-level outcomes, and missed releases give no consistent reliability
ranking. Execution granularity should therefore be chosen against measured
complete-path headroom and age requirements; broader platform and
audio-device claims require further measurement.
